\documentclass[11pt]{article}
\usepackage{amsmath,amssymb,amsthm}
\usepackage[margin=2.5cm]{geometry}

\newtheorem*{lemmathree}{Lemma 3}
\newtheorem*{lemmasix}{Lemma 6}
\newtheorem*{theoremnine}{Theorem 9}

\title{Problem B (Bulgarian solitaire), Cell C1: Cyclic Partitions and Cycles\\[4pt]
\large Written proof}
\author{Submission of team record B-C1}
\date{}

\begin{document}
\maketitle

\begin{center}
\fbox{\begin{minipage}{0.85\textwidth}
\textbf{Cell status:} SOLVED\qquad \textbf{Claim status:} PROVED\\[3pt]
Established: parts (i), (ii)(a) and (ii)(b) of Cell C1, for every $k\ge1$ and every $1\le r\le k$.
The proof is purely written; the included script (Section~\ref{sec:verify}) is a sanity check for
$1\le n\le 45$ and the proof does not rely on it.
\end{minipage}}
\end{center}

\section{Statement}

\subsection*{Definitions (verbatim from the official problem statement)}

A \emph{partition} of a positive integer $n$ is a weakly decreasing sequence
\[
\lambda=(\lambda_1,\ldots,\lambda_s)
\]
of positive integers with
\[
\lambda_1+\cdots+\lambda_s=n.
\]
Think of it as a division of $n$ cards into $s$ piles.

The \emph{shift} $B(\lambda)$ is the partition of $n$ obtained as follows: remove one card from every pile,
discard the piles that become empty, and add one new pile consisting of the $s$ removed cards. Formally,
$B(\lambda)$ is the partition whose parts are the positive numbers among
\[
\lambda_1-1,\;\ldots,\;\lambda_s-1
\]
together with one extra part equal to $s$.

Call $\lambda$ \emph{cyclic} if
\[
B^i(\lambda)=\lambda\qquad\text{for some } i\ge 1.
\]

Write
\[
T_k=\frac{k(k+1)}{2},\qquad \delta_k=(k,k-1,\ldots,2,1),
\]
the partition of $T_k$ into distinct parts. Every $n\ge 1$ satisfies
\[
T_{k-1}<n\le T_k
\]
for exactly one $k$, which we call the \emph{rank} of $n$.

\subsection*{Cell C1 (verbatim)}

\emph{First, the long-run behaviour: which partitions repeat under the shift, and how they fall into cycles.
Let $n=T_k$. Prove that for every partition $\lambda$ of $n$ there is an $i$ with $B^i(\lambda)=\delta_k$, and
that $\delta_k$ is the only cyclic partition of $n$. Then let $n$ be arbitrary of rank $k$, say
$n=T_{k-1}+r$ with $1\le r\le k$: determine all cyclic partitions of $n$, and determine the number of
distinct cycles of $B$ on the partitions of $n$. Prove both.}

\subsection*{Exact targets (as fixed in the team's target file)}

\begin{itemize}
\item[(i)] For every $k\ge1$ and $n=T_k$: for every partition $\lambda$ of $n$ there is an $i\ge0$ with
$B^i(\lambda)=\delta_k$; and $\delta_k$ is the only cyclic partition of $n$.
\item[(ii)] For every $k\ge1$ and every $r$ with $1\le r\le k$, $n=T_{k-1}+r$:
\begin{itemize}
\item[(a)] give an explicit description (as functions of $k$ and $r$) of the set of \emph{all} cyclic
partitions of $n$, and prove that a partition of $n$ is cyclic if and only if it lies in that set;
\item[(b)] give the exact number of distinct cycles of $B$ on the partitions of $n$ (as a formula in $k$
and $r$), and prove it.
\end{itemize}
\end{itemize}

\section{Status}

\begin{itemize}
\item \textbf{Cell status:} SOLVED.
\item \textbf{Claim status:} PROVED.
\end{itemize}

\section{Answer}

Let $k\ge1$, $1\le r\le k$, $n=T_{k-1}+r$.
\begin{itemize}
\item[(ii)(a)] A partition of $n$ is cyclic if and only if it equals
\[
\lambda(\varepsilon):=(k-1+\varepsilon_1,\,k-2+\varepsilon_2,\,\ldots,\,1+\varepsilon_{k-1},\,\varepsilon_k)
\quad\text{(last entry deleted when it is $0$)}
\]
for some $\varepsilon\in\{0,1\}^k$ with $\varepsilon_1+\cdots+\varepsilon_k=r$. Equivalently: the Young diagram
is the staircase $\delta_{k-1}$ plus exactly $r$ of the $k$ cells of the $k$-th diagonal. There are exactly
$\binom kr$ cyclic partitions of $n$, and $B$ acts on them by
$B(\lambda(\varepsilon))=\lambda(\varepsilon_k,\varepsilon_1,\ldots,\varepsilon_{k-1})$.
\item[(ii)(b)] The number of distinct cycles of $B$ on the partitions of $n$ is
\[
N(k,r)=\frac1k\sum_{d\mid\gcd(k,r)}\varphi(d)\binom{k/d}{r/d}.
\]
\item[(i)] (Case $r=k$, $n=T_k$.) $\delta_k$ is the only cyclic partition of $T_k$, and every partition
$\lambda$ of $T_k$ has $B^i(\lambda)=\delta_k$ for some $i\ge0$.
\end{itemize}

\section{Cited results vs.\ our contribution}

\paragraph{Cited (standard results used by name in the proof).}
\begin{itemize}
\item Pigeonhole principle (step 1.2).
\item Chinese remainder theorem for the coprime moduli $d$, $d+1$ (proof of Lemma 6).
\item Orbit-counting lemma (Burnside's lemma), step 10.2; the proof includes a short derivation of it
from the orbit--stabiliser theorem, which is used without proof.
\item Euler's totient function $\varphi$ (convention 0.1); the identity $\sum_{d\mid k}\varphi(d)=k$ used in the
check after 10.4 is justified there by partitioning $\{1,\ldots,k\}$ by $\gcd(\cdot,k)$.
\end{itemize}
The accepted artefacts give no bibliographic references for these standard facts, and none are added here.
The accepted artefacts cite no literature on Bulgarian solitaire, and no literature search is part of the
material this document was prepared from; this document therefore makes no claim either way about whether
these results appear in the literature.

\paragraph{Our contribution (the team's accepted artefacts).}
\begin{itemize}
\item The written proof below (Sections 0--11 of the proof), establishing claims R1--R11, all with claim
status PROVED.
\item The sanity-check script \texttt{code/check\_c1.py} (stdlib only, exact integer arithmetic), which confirms
(i), (ii)(a), (ii)(b) by brute force for all $1\le n\le 45$. It is not load-bearing.
\end{itemize}

\section{Proof}

The text of this section is the team's accepted proof, typeset from Markdown to \LaTeX{} with its numbering
preserved. The only editorial change is the script path in the final paragraph (see Section~\ref{sec:lim}).

\paragraph{Statement proved.} For every $k\ge 1$ and every $1\le r\le k$, with $n=T_{k-1}+r$: a partition of $n$
is cyclic if and only if it equals
$\lambda(\varepsilon):=(k-1+\varepsilon_1,\,k-2+\varepsilon_2,\,\ldots,\,1+\varepsilon_{k-1},\,\varepsilon_k)$
(last entry deleted when it is $0$) for some $\varepsilon\in\{0,1\}^k$ with $\varepsilon_1+\cdots+\varepsilon_k=r$;
$B$ acts on these by $B(\lambda(\varepsilon))=\lambda(\varepsilon_k,\varepsilon_1,\ldots,\varepsilon_{k-1})$; the
number of distinct cycles of $B$ on the partitions of $n$ is
$N(k,r)=\frac1k\sum_{d\mid\gcd(k,r)}\varphi(d)\binom{k/d}{r/d}$. In particular (case $r=k$, $n=T_k$) $\delta_k$ is
the only cyclic partition of $T_k$ and every partition $\lambda$ of $T_k$ has $B^i(\lambda)=\delta_k$ for some
$i\ge 0$.

\medskip\hrule\medskip

\subsection*{0. Conventions}

\paragraph{0.1.} Reading of ``cycle of $B$'': a \emph{cycle} is a set $\{\mu,B(\mu),\ldots,B^{p-1}(\mu)\}$ where
$B^p(\mu)=\mu$, $p\ge1$ (the forward orbit of a cyclic partition). Two cycles are distinct if they are distinct
as sets. $\varphi$ is Euler's totient function. $\mathbb N_{\ge1}=\{1,2,3,\ldots\}$.

\paragraph{0.2.} For a finite sequence $c=(c_1,\ldots,c_m)$ of positive integers (not necessarily ordered) put
\[
\mathcal C(c)=\{(j,h)\in\mathbb N_{\ge1}^2:\ 1\le j\le m,\ 1\le h\le c_j\}
\]
(pile $j$, height $h$), so $|\mathcal C(c)|=c_1+\cdots+c_m$. A partition is determined by its cell set:
$\lambda_j=|\{h:(j,h)\in\mathcal C(\lambda)\}|$.

\paragraph{0.3.} For a partition $\nu$: if $(1,h)\in\mathcal C(\nu)$ and $1\le h'\le h$ then
$(1,h')\in\mathcal C(\nu)$ (definition of $\mathcal C$ with $j=1$).

\paragraph{0.4.} Diagonals: for $d\ge1$ let
$D_d=\{(j,h)\in\mathbb N_{\ge1}^2: j+h-1=d\}=\{(j,d+1-j):1\le j\le d\}$, so $|D_d|=d$, and every cell of
$\mathbb N_{\ge1}^2$ lies in exactly one $D_d$ (namely $d=j+h-1\ge1$).

\paragraph{0.5.} Weight: $E(c)=\sum_{(j,h)\in\mathcal C(c)}(j+h-1)$. Summing over $h$ first,
\[
E(c)=\sum_{j=1}^m\sum_{h=1}^{c_j}(j-1+h)=\sum_{j=1}^m (j-1)c_j+\sum_{j=1}^m\frac{c_j(c_j+1)}2 .
\]

\paragraph{0.6.} Unsorted shift: for a partition $\lambda=(\lambda_1,\ldots,\lambda_s)$ of $n$ let
$t=|\{j:\lambda_j\ge2\}|$. Because $\lambda$ is weakly decreasing, $\lambda_j\ge2\iff j\le t$. Put
\[
U(\lambda)=(s,\ \lambda_1-1,\ \ldots,\ \lambda_t-1),
\]
a sequence of $t+1$ positive integers. Its entries are exactly the positive numbers among
$\lambda_1-1,\ldots,\lambda_s-1$ together with $s$, so by the definition of the shift, $B(\lambda)$ is the weakly
decreasing rearrangement of $U(\lambda)$; write $B(\lambda)=\operatorname{sort}(U(\lambda))$.

\paragraph{0.7.} The map $R:\mathbb N_{\ge1}^2\to\mathbb N_{\ge1}^2$: $R(j,h)=(j+1,h-1)$ if $h\ge2$, and
$R(j,1)=(1,j)$.

\subsection*{1. $B$ is a self-map; orbits reach cyclic partitions (R1)}

\paragraph{1.1.} $B(\lambda)$ is a partition of $n$: its parts are positive (the $\lambda_j-1$ kept are positive
and $s\ge1$), it is weakly decreasing by 0.6, and its sum is
$\sum_{j\le t}(\lambda_j-1)+s=\sum_{j\le s}(\lambda_j-1)+s=n$, since $\lambda_j-1=0$ for $j>t$.

\paragraph{1.2.} The set $P(n)$ of partitions of $n$ is finite (each is a sequence of at most $n$ entries from
$\{1,\ldots,n\}$). For $\lambda\in P(n)$, the elements $B^0(\lambda),B^1(\lambda),\ldots,B^{|P(n)|}(\lambda)$ are
$|P(n)|+1$ elements of $P(n)$, so by the pigeonhole principle $B^a(\lambda)=B^b(\lambda)$ for some $0\le a<b$.
Then $\mu=B^a(\lambda)$ satisfies $B^{b-a}(\mu)=\mu$ with $b-a\ge1$, i.e.\ $\mu$ is cyclic.

\subsection*{2. The unsorted shift moves cells by $R$ (R2)}

\paragraph{2.1.} Claim: $R$ restricts to a bijection $\mathcal C(\lambda)\to\mathcal C(U(\lambda))$. Write
$U(\lambda)=(c_1,\ldots,c_{t+1})$ with $c_1=s$, $c_{j+1}=\lambda_j-1$ ($1\le j\le t$).
\begin{itemize}
\item If $(j,h)\in\mathcal C(\lambda)$ with $h\ge2$: then $\lambda_j\ge h\ge2$, so $j\le t$ (0.6), and
$R(j,h)=(j+1,h-1)$ has $2\le j+1\le t+1$ and $1\le h-1\le\lambda_j-1=c_{j+1}$; so
$R(j,h)\in\mathcal C(U(\lambda))$ with first coordinate $\ge2$.
\item If $(j,1)\in\mathcal C(\lambda)$: $1\le j\le s=c_1$, so $R(j,1)=(1,j)\in\mathcal C(U(\lambda))$ with first
coordinate $1$.
\item Inverse map $\mathcal C(U(\lambda))\to\mathcal C(\lambda)$: $(1,h)\mapsto(h,1)$ (valid: $h\le c_1=s$, and
$\lambda_h\ge1$); $(j',h')\mapsto(j'-1,h'+1)$ for $j'\ge2$ (valid: $j'-1\le t$ and
$h'+1\le c_{j'}+1=\lambda_{j'-1}$). Composing in both orders gives the identity (check on the two cases $h=1$ /
$h\ge2$, respectively $j'=1$ / $j'\ge2$). So $R:\mathcal C(\lambda)\to\mathcal C(U(\lambda))$ is a bijection.
\end{itemize}

\paragraph{2.2.} $R$ preserves $j+h-1$: $(j+1)+(h-1)-1=j+h-1$ and $1+j-1=j+1-1$. Hence $R(D_d)\subseteq D_d$ for
each $d$, and by 2.1, $E(U(\lambda))=E(\lambda)$.

\subsection*{3. Sorting lemma (R3)}

\begin{lemmathree}
Let $c=(c_1,\ldots,c_m)$ be positive integers with sum $n$ and let $\lambda=\operatorname{sort}(c)$ be its
weakly decreasing rearrangement. Then $E(\lambda)\le E(c)$, with equality if and only if $c$ is weakly
decreasing (i.e.\ $c=\lambda$).
\end{lemmathree}

\paragraph{3.1.} By 0.5 and since $\sum_j c_j(c_j+1)/2$ is unchanged by rearranging,
$E(c)-E(\lambda)=\sum_j(j-1)c_j-\sum_j(j-1)\lambda_j$.

\paragraph{3.2.} Exchanging the order of summation,
$\sum_{j=1}^m (j-1)c_j=\sum_{j=1}^m\sum_{i=1}^{j-1}c_j=\sum_{i=1}^{m-1}\sum_{j=i+1}^m c_j=\sum_{i=1}^{m-1}\bigl(n-P_i(c)\bigr)$
with $P_i(c)=c_1+\cdots+c_i$. The same identity holds for $\lambda$.

\paragraph{3.3.} For each $i$, $P_i(c)\le P_i(\lambda)$: there is a permutation $\sigma$ with
$c_l=\lambda_{\sigma(l)}$; list $\sigma(1),\ldots,\sigma(i)$ increasingly as $j_1<\cdots<j_i$; then $j_l\ge l$, so
$\lambda_{j_l}\le\lambda_l$ ($\lambda$ weakly decreasing), and summing,
$P_i(c)=\sum_l\lambda_{j_l}\le\sum_{l\le i}\lambda_l=P_i(\lambda)$.

\paragraph{3.4.} By 3.2 and 3.3, $E(c)-E(\lambda)=\sum_{i=1}^{m-1}(P_i(\lambda)-P_i(c))\ge0$.

\paragraph{3.5.} Equality case. If $c$ is weakly decreasing then $c=\lambda$ (the weakly decreasing rearrangement
of a multiset is unique) and equality holds. If $c$ is not weakly decreasing, pick $i$ with $c_i<c_{i+1}$
($1\le i\le m-1$). The $i$ entries $c_1,\ldots,c_{i-1},c_{i+1}$ (at distinct positions) have sum
$P_i(c)-c_i+c_{i+1}>P_i(c)$, and by the argument of 3.3 applied to these $i$ positions, that sum is
$\le P_i(\lambda)$. So $P_i(c)<P_i(\lambda)$, and by 3.4 $E(c)>E(\lambda)$. $\square$

\subsection*{4. Energy along orbits; no sorting on cycles (R4)}

\paragraph{4.1.} For every partition $\lambda$:
$E(B(\lambda))=E(\operatorname{sort}(U(\lambda)))\le E(U(\lambda))=E(\lambda)$ (Lemma 3, then 2.2), with equality
iff $U(\lambda)$ is weakly decreasing, in which case $B(\lambda)=U(\lambda)$ as sequences and so, by 2.1,
$\mathcal C(B(\lambda))=R(\mathcal C(\lambda))$.

\paragraph{4.2.} Let $\lambda$ be cyclic, $B^p(\lambda)=\lambda$, $p\ge1$. By 4.1,
$E(\lambda)\ge E(B\lambda)\ge\cdots\ge E(B^p\lambda)=E(\lambda)$, so $E(B^{i+1}\lambda)=E(B^i\lambda)$ for
$0\le i<p$. For arbitrary $i\ge0$, $B^i\lambda=B^{i\bmod p}\lambda$ (induction on $i$ using $B^p\lambda=\lambda$),
so $E(B^{i+1}\lambda)=E(B^i\lambda)$ for all $i\ge0$. By the equality case of 4.1,
$\mathcal C(B^{i+1}\lambda)=R(\mathcal C(B^i\lambda))$ for all $i\ge0$, hence by induction
\[
\mathcal C(B^t\lambda)=R^t(\mathcal C(\lambda))\qquad(t\ge0).
\]

\subsection*{5. $R$ rotates each diagonal (R5)}

\paragraph{5.1.} Index $D_d$ by its first coordinate: $x_a=(a,d+1-a)$, $1\le a\le d$. For $a<d$ the height
$d+1-a\ge2$, so $R(x_a)=(a+1,d-a)=x_{a+1}$; for $a=d$, $R(x_d)=R(d,1)=(1,d)=x_1$. So $R|_{D_d}$ is the cyclic
rotation $x_a\mapsto x_{a+1}$ (indices mod $d$ in $\{1,\ldots,d\}$), a bijection of $D_d$, and by induction on
$t$, $R^t(x_a)=x_{a'}$ with $a'\equiv a+t\pmod d$.

\paragraph{5.2.} Consequently, for cyclic $\lambda$ and $t\ge0$, by 4.2 and 5.1 (each $R^t$ maps $D_d$
bijectively onto itself and maps cells in different diagonals to different diagonals):
\[
\mathcal C(B^t\lambda)\cap D_d=R^t\bigl(\mathcal C(\lambda)\cap D_d\bigr),
\]
and for $x\in D_d$: $x\in\mathcal C(\lambda)\iff R^t(x)\in\mathcal C(B^t\lambda)$ (injectivity of $R^t$ on $D_d$).

\subsection*{6. Key lemma (R6)}

\begin{lemmasix}
Let $\lambda$ be cyclic and $d\ge1$. If $D_d\not\subseteq\mathcal C(\lambda)$ then
$D_{d+1}\cap\mathcal C(\lambda)=\emptyset$.
\end{lemmasix}

\begin{proof}
Suppose $x_a=(a,d+1-a)\in D_d\setminus\mathcal C(\lambda)$ and $y_b=(b,d+2-b)\in D_{d+1}\cap\mathcal C(\lambda)$.
Since $\gcd(d,d+1)=1$, the Chinese remainder theorem gives $t\ge0$ with $t\equiv1-a\pmod d$ and
$t\equiv1-b\pmod{d+1}$. By 5.1, $R^t(x_a)=(1,d)$ and $R^t(y_b)=(1,d+1)$. By 5.2,
$(1,d)\notin\mathcal C(B^t\lambda)$ and $(1,d+1)\in\mathcal C(B^t\lambda)$. But $B^t\lambda$ is a partition (1.1),
and by 0.3, $(1,d+1)\in\mathcal C(B^t\lambda)$ forces $(1,d)\in\mathcal C(B^t\lambda)$. Contradiction.
\end{proof}

\subsection*{7. Structure of cyclic partitions (R7)}

\paragraph{7.1.} Let $\lambda$ be a cyclic partition of $n$. $\mathcal C(\lambda)$ is finite and $D_d$ is
nonempty, so $K=\min\{d\ge1: D_d\not\subseteq\mathcal C(\lambda)\}$ exists. $D_1=\{(1,1)\}\subseteq\mathcal C(\lambda)$
since $\lambda_1\ge1$, so $K\ge2$.

\paragraph{7.2.} Claim: $D_e\cap\mathcal C(\lambda)=\emptyset$ for all $e>K$. Induction on $e$. For $e=K+1$:
Lemma 6 with $d=K$. For $e>K+1$: by induction $D_{e-1}\cap\mathcal C(\lambda)=\emptyset$; since
$D_{e-1}\neq\emptyset$, $D_{e-1}\not\subseteq\mathcal C(\lambda)$, and Lemma 6 with $d=e-1$ gives
$D_e\cap\mathcal C(\lambda)=\emptyset$.

\paragraph{7.3.} By 0.4, 7.1, 7.2: $\mathcal C(\lambda)=D_1\cup\cdots\cup D_{K-1}\cup S$ with
$S=\mathcal C(\lambda)\cap D_K\subsetneq D_K$, so $n=|\mathcal C(\lambda)|=T_{K-1}+|S|$ with $0\le|S|\le K-1$.

\paragraph{7.4.} Rank. $T_0<T_1<T_2<\cdots$ (as $T_m-T_{m-1}=m\ge1$) and $T_0=0$, so the intervals
$(T_{m-1},T_m]$, $m\ge1$, are disjoint and cover $\mathbb N_{\ge1}$: the rank is well defined. Let
$n=T_{k-1}+r$, $1\le r\le k$, i.e.\ $n$ has rank $k$.
\begin{itemize}
\item If $|S|\ge1$: $T_{K-1}<n\le T_{K-1}+K-1<T_K$, so the rank is $K$: $k=K$ and $r=n-T_{k-1}=|S|$.
\item If $|S|=0$: $n=T_{K-1}$ with $K-1\ge1$, and $T_{K-2}<n\le T_{K-1}$, so $k=K-1$ and
$r=n-T_{k-1}=T_k-T_{k-1}=k$; here $\mathcal C(\lambda)=D_1\cup\cdots\cup D_k$.
\end{itemize}
In both cases:
\begin{equation*}
D_1\cup\cdots\cup D_{k-1}\subseteq\mathcal C(\lambda)\subseteq D_1\cup\cdots\cup D_k,\qquad
|\mathcal C(\lambda)\cap D_k|=r.\tag{$*$}
\end{equation*}

\paragraph{7.5.} Reading off the parts. Define $\varepsilon_j=1$ if $(j,k+1-j)\in\mathcal C(\lambda)$ and
$\varepsilon_j=0$ otherwise ($1\le j\le k$); by $(*)$, $\sum_j\varepsilon_j=r$. For $1\le j\le k$ the cells $(j,h)$
with $h\le k-j$ lie in $D_{j+h-1}$ with $j+h-1\le k-1$, so are in $\mathcal C(\lambda)$; the cell
$(j,k+1-j)\in D_k$ is present iff $\varepsilon_j=1$; cells $(j,h)$ with $h\ge k+2-j$ lie in $D_e$, $e\ge k+1$, so
are absent. Hence $\lambda_j=k-j+\varepsilon_j$ for $1\le j\le k$ (where $\lambda_j:=0$ if $j>s$). For $j>k$, every
cell $(j,h)$ lies in $D_e$ with $e\ge j\ge k+1$, so $\lambda_j=0$. Therefore $\lambda=\lambda(\varepsilon)$ as
defined in the statement line.

\subsection*{8. The converse and the action of $B$ (R8)}

\paragraph{8.1.} Let $W(k,r)=\{\varepsilon\in\{0,1\}^k:\sum_j\varepsilon_j=r\}$ and for $\varepsilon\in W(k,r)$ let
$\lambda(\varepsilon)$ be the sequence of positive numbers among $k-j+\varepsilon_j$, $j=1,\ldots,k$, in this order.
For $j\le k-1$, $k-j+\varepsilon_j\ge1$, and the $k$-th value is $\varepsilon_k$, so only the last entry can be
dropped (exactly when $\varepsilon_k=0$). The sequence is weakly decreasing since
$(k-j+\varepsilon_j)-(k-j-1+\varepsilon_{j+1})=1+\varepsilon_j-\varepsilon_{j+1}\ge0$, and its sum is
$\sum_{j=1}^k(k-j)+r=T_{k-1}+r=n$. So $\lambda(\varepsilon)$ is a partition of $n$.

\paragraph{8.2.} Injectivity: from $\lambda=\lambda(\varepsilon)$ (with $\lambda_j:=0$ beyond its length) we recover
$\varepsilon_j=\lambda_j-(k-j)$, $1\le j\le k$.

\paragraph{8.3.} Let $\rho:W(k,r)\to W(k,r)$, $\rho(\varepsilon)=(\varepsilon_k,\varepsilon_1,\ldots,\varepsilon_{k-1})$.
Claim: $B(\lambda(\varepsilon))=\lambda(\rho\varepsilon)$. Write $\lambda=\lambda(\varepsilon)$, with
$s=k-1+\varepsilon_k$ parts. The parts of $B(\lambda)$ are $s=k-1+\varepsilon_k$ together with the positive values
among $\lambda_j-1$: for $1\le j\le k-1$, $\lambda_j-1=k-1-j+\varepsilon_j$; if $\varepsilon_k=1$ there is also
$\lambda_k-1=0$, which is not positive. So the multiset of parts of $B(\lambda)$ is
$\{k-1+\varepsilon_k\}\cup\{k-(j+1)+\varepsilon_j:1\le j\le k-1,\ \text{value}>0\}$. With $\varepsilon'=\rho\varepsilon$
($\varepsilon'_1=\varepsilon_k$, $\varepsilon'_{j+1}=\varepsilon_j$), the parts of $\lambda(\varepsilon')$ are the
positive values among $k-1+\varepsilon'_1=k-1+\varepsilon_k$ (positive, as $k-1+\varepsilon_k=s\ge1$) and
$k-(j+1)+\varepsilon'_{j+1}=k-(j+1)+\varepsilon_j$, $1\le j\le k-1$: the same multiset. A partition is determined by
its multiset of parts, so $B(\lambda(\varepsilon))=\lambda(\rho\varepsilon)$. (For $k=1$: $W(1,1)=\{(1)\}$,
$\lambda=(1)$, $B((1))=(1)$, consistent.)

\paragraph{8.4.} By 8.3 and induction, $B^i(\lambda(\varepsilon))=\lambda(\rho^i\varepsilon)$; since
$\rho^k=\mathrm{id}$ (each application moves every letter one place cyclically),
$B^k(\lambda(\varepsilon))=\lambda(\varepsilon)$ with $k\ge1$, so $\lambda(\varepsilon)$ is cyclic.

\subsection*{9. Answer to (ii)(a) (R9)}

\begin{theoremnine}
For $k\ge1$, $1\le r\le k$, $n=T_{k-1}+r$, the cyclic partitions of $n$ are exactly the partitions
$\lambda(\varepsilon)$, $\varepsilon\in W(k,r)$; equivalently, the partitions whose Young diagram consists of the
staircase $\delta_{k-1}$ (all cells on diagonals $D_1,\ldots,D_{k-1}$) plus exactly $r$ of the $k$ cells of
diagonal $D_k$. There are exactly $\binom kr$ of them.
\end{theoremnine}

\begin{proof}
``Only if'': 7.5. ``If'': 8.4. The count follows from the injectivity 8.2 and $|W(k,r)|=\binom kr$.
\end{proof}

\subsection*{10. Answer to (ii)(b) (R10)}

\paragraph{10.1.} Let $\mathrm{Cyc}(n)$ be the set of cyclic partitions and
$\Lambda:W(k,r)\to\mathrm{Cyc}(n)$, $\varepsilon\mapsto\lambda(\varepsilon)$: a bijection by Theorem 9 and 8.2,
with $B\circ\Lambda=\Lambda\circ\rho$ (8.3). Every cycle (0.1) consists of cyclic partitions: if $B^p\mu=\mu$
then $B^p(B^i\mu)=B^i(B^p\mu)=B^i\mu$. For cyclic $\mu=\Lambda(\varepsilon)$, the cycle through $\mu$ is
$\{B^i\mu:i\ge0\}=\Lambda(\{\rho^i\varepsilon:0\le i<k\})$, the image of the orbit of $\varepsilon$ under the cyclic
group $G=\langle\rho\rangle$ (of order dividing $k$; we let $G$ act through $j\mapsto\rho^j$, $j\in\mathbb Z/k$).
Two cyclic partitions lie in the same cycle iff their $\varepsilon$'s lie in the same orbit, and orbits partition
$W(k,r)$. Hence the number of distinct cycles equals the number of orbits of $\mathbb Z/k$ acting on $W(k,r)$ by
$j\cdot\varepsilon=\rho^j\varepsilon$.

\paragraph{10.2.} Orbit-counting lemma (Burnside's lemma, standard; proof included): for a finite group $H$ acting
on a finite set $X$, $\#\text{orbits}=\frac1{|H|}\sum_{g\in H}|\mathrm{Fix}(g)|$. Proof:
$\sum_g|\mathrm{Fix}(g)|=|\{(g,x):gx=x\}|=\sum_x|\mathrm{Stab}(x)|=\sum_x|H|/|Hx|$ (orbit--stabiliser theorem), and
$\sum_{x\in O}1/|O|=1$ for each orbit $O$.

\paragraph{10.3.} Fixed points. Index positions by $\mathbb Z/k$; $(\rho^j\varepsilon)_i=\varepsilon_{i-j}$. So
$\rho^j\varepsilon=\varepsilon$ iff $\varepsilon_i=\varepsilon_{i-j}$ for all $i$, iff $\varepsilon$ is constant on each
coset of the subgroup $\langle j\rangle\le\mathbb Z/k$. With $g=\gcd(j,k)$ (and $\gcd(0,k)=k$),
$\langle j\rangle=g\mathbb Z/k\mathbb Z$ has order $k/g$ and there are $g$ cosets. A word constant on cosets is
determined by the set of cosets where it equals $1$; it has $r$ ones iff that set has $r/(k/g)$ elements. Hence
$|\mathrm{Fix}(\rho^j)|=\binom{g}{rg/k}$ if $(k/g)\mid r$, and $0$ otherwise.

\paragraph{10.4.} For each divisor $g$ of $k$, the $j\in\{0,\ldots,k-1\}$ with $\gcd(j,k)=g$ are $j=gj'$ with
$0\le j'<k/g$ and $\gcd(j',k/g)=1$; there are $\varphi(k/g)$ of them. Substituting $d=k/g$ (so $d\mid k$ and the
condition $(k/g)\mid r$ reads $d\mid r$):
\[
\#\text{cycles}=N(k,r)=\frac1k\sum_{d\mid\gcd(k,r)}\varphi(d)\binom{k/d}{r/d}.
\]
Checks: $r=k$ gives $\frac1k\sum_{d\mid k}\varphi(d)=1$ (by $\sum_{d\mid k}\varphi(d)=k$: partition
$\{1,\ldots,k\}$ by $\gcd(\cdot,k)$ as in 10.4); $\gcd(k,r)=1$ gives $\binom kr/k$. $\square$

\paragraph{10.5.} (Remark, not needed.) The cycle through $\lambda(\varepsilon)$ has length equal to the least
$p\ge1$ with $\rho^p\varepsilon=\varepsilon$ (by 8.2--8.3), a divisor of $k$.

\subsection*{11. Part (i) (R11)}

Let $n=T_k$, i.e.\ $r=k$ in Theorem 9 (rank of $T_k$ is $k$ by 7.4). $W(k,k)=\{(1,\ldots,1)\}$ and
$\lambda(1,\ldots,1)=(k,k-1,\ldots,1)=\delta_k$, so $\delta_k$ is the only cyclic partition of $T_k$ (and
$B(\delta_k)=\delta_k$ by 8.3). For any partition $\lambda$ of $T_k$, 1.2 gives $a\ge0$ with $B^a(\lambda)$
cyclic, hence $B^a(\lambda)=\delta_k$. $\square$

\subsection*{Established / not established}

Established: (i) and (ii)(a),(b) for all $k\ge1$, $1\le r\le k$, with the answer above. The brute-force script
\texttt{code/check\_c1.py} (all $n\le45$) agrees but the proof does not rely on it. Not addressed: time-to-cycle
questions (cells C2 and beyond). Known gaps: none found.

\section{How to verify}\label{sec:verify}

\paragraph{The proof.} The proof above is self-contained and is to be checked by reading; it does not rely on any
computation.

\paragraph{Sanity check (not load-bearing).} The script \texttt{code/check\_c1.py} (listed in full in the
appendix; Python standard library only, exact integer arithmetic) does the following for every $n$ in
$[1,\mathrm{NMAX}]$: enumerates all partitions of $n$; finds the cyclic ones by iterating $B$; checks that they
equal the predicted set $\{\lambda(\varepsilon):\varepsilon\in\{0,1\}^k,\ |\varepsilon|=r\}$ ($k$ the rank,
$r=n-T_{k-1}$); counts the cycles and compares with $\frac1k\sum_{d\mid\gcd(k,r)}\varphi(d)\binom{k/d}{r/d}$;
and for triangular $n$ checks that every orbit reaches $\delta_k$. It checks nothing beyond $n\le\mathrm{NMAX}$.

Run from the submission folder (the folder containing \texttt{code/}):
\begin{verbatim}
/usr/bin/time -p python3 code/check_c1.py 45
\end{verbatim}
Expected output: 46 lines, one line per $n=1,\ldots,45$ of the form
\begin{verbatim}
n=45 k=9 r=9 #partitions=89134 #cyclic=1 #cycles=1 OK
\end{verbatim}
followed by the final line
\begin{verbatim}
ALL OK up to 45
\end{verbatim}
Measured runtime (Python 3.14.0, this submission's copy of the script): real 19.11\,s with the project
virtual-environment interpreter and real 20.86\,s with the system \texttt{python3}; both runs exited with
status 0 and produced byte-identical output. This is well under the 10-minute limit.

\section{Limitations}\label{sec:lim}

\begin{itemize}
\item What is established: (i), (ii)(a) and (ii)(b) of Cell C1, for all $k\ge1$ and $1\le r\le k$. Not
addressed: the time-to-cycle questions (cells C2 and beyond).
\item The count in (ii)(b) is for the reading of ``distinct cycles'' fixed in convention 0.1: a cycle is the
forward orbit of a cyclic partition, and two cycles are distinct if they are distinct as sets.
\item The computer check covers only $1\le n\le 45$ and proves nothing beyond that range; the proof does not
rely on it.
\item Standard results (pigeonhole principle, Chinese remainder theorem, orbit--stabiliser theorem) are used
without proof and without bibliographic reference; see the Cited section.
\item No literature search is part of the material this document was prepared from; no claim is made about
whether these results appear in the literature.
\item Editorial change: in the final paragraph of the proof (``Established / not established''), the accepted
proof names the script as \texttt{out/code/check\_c1.py}; in this submission the script is at
\texttt{code/check\_c1.py}, and the path has been changed accordingly. No other change was made to the text of
the proof beyond conversion from Markdown to \LaTeX.
\item This document was typeset by the team's scribe from the accepted artefacts. The scribe added no
mathematics and did not referee the proof.
\end{itemize}

\appendix
\section{Listing of \texttt{code/check\_c1.py}}

\begingroup\footnotesize
\begin{verbatim}
"""Sanity check for B-C1 (stdlib only, exact integer arithmetic).
For every n in [1, NMAX]: enumerate all partitions of n, compute the cyclic ones by
iterating B (functional graph), compare with the predicted set {lambda(eps)}, count
cycles and compare with (1/k) sum_{d | gcd(k,r)} phi(d) C(k/d, r/d).
Also checks for triangular n that every orbit reaches delta_k.
Not load-bearing: the proof does not rest on this computation."""
import sys
from math import comb, gcd

def partitions(n, maxpart=None):
    if maxpart is None:
        maxpart = n
    if n == 0:
        yield ()
        return
    for p in range(min(n, maxpart), 0, -1):
        for rest in partitions(n - p, p):
            yield (p,) + rest

def B(lam):
    s = len(lam)
    parts = [x - 1 for x in lam if x > 1] + [s]
    return tuple(sorted(parts, reverse=True))

def phi(m):
    return sum(1 for i in range(1, m + 1) if gcd(i, m) == 1)

def rank(n):
    k = 1
    while k * (k + 1) // 2 < n:
        k += 1
    return k

def predicted(k, r):
    out = set()
    for mask in range(1 << k):
        eps = [(mask >> (j)) & 1 for j in range(k)]  # eps[0..k-1] = eps_1..eps_k
        if sum(eps) != r:
            continue
        lam = tuple(x for x in (k - (j + 1) + eps[j] for j in range(k)) if x > 0)
        out.add(lam)
    return out

def main(NMAX):
    for n in range(1, NMAX + 1):
        k = rank(n)
        r = n - (k - 1) * k // 2
        parts = list(partitions(n))
        cyc = set()
        for lam in parts:
            seen = {}
            x = lam
            i = 0
            while x not in seen:
                seen[x] = i
                x = B(x)
                i += 1
            # x is the first repeated element: it lies on the cycle
            y = x
            while True:
                cyc.add(y)
                y = B(y)
                if y == x:
                    break
        pred = predicted(k, r)
        assert cyc == pred, (n, k, r)
        # count cycles
        left = set(cyc)
        ncyc = 0
        while left:
            x = left.pop()
            y = B(x)
            while y != x:
                left.discard(y)
                y = B(y)
            ncyc += 1
        g = gcd(k, r)
        num = sum(phi(d) * comb(k // d, r // d) for d in range(1, g + 1) if g % d == 0)
        assert num % k == 0 and num // k == ncyc, (n, k, r, ncyc, num)
        if r == k:
            delta = tuple(range(k, 0, -1))
            assert cyc == {delta}
            for lam in parts:
                x = lam
                for _ in range(len(parts) + 1):
                    if x == delta:
                        break
                    x = B(x)
                assert x == delta, (n, lam)
        print(f"n={n} k={k} r={r} #partitions={len(parts)} #cyclic={len(cyc)} #cycles={ncyc} OK")
    print("ALL OK up to", NMAX)

if __name__ == "__main__":
    main(int(sys.argv[1]) if len(sys.argv) > 1 else 40)
\end{verbatim}
\endgroup

\end{document}
